% Level 1, Week 5 -- workbook chapter
\chapterweek{5}{Messages and My Blocks}{Break big programs into named parts.}

\section{What you will learn}
\begin{itemize}[leftmargin=1.2em,itemsep=1pt]
  \item How to break a big job into smaller named parts (\textbf{decomposition}).
  \item How to make your own block with \textbf{My Blocks}, and give it an \textbf{input}.
  \item How sprites send each other \textbf{messages} with \textbf{broadcast}.
  \item How to pick colors and text that \textbf{everyone} can see.
\end{itemize}

\section{The big idea}
\subsection{Big jobs, small named parts}
A garden picture is a big job. Break it into parts and give each part a name. Then the main plan is short and easy to read.

\begin{center}
\begin{tikzpicture}[level distance=1.5cm,sibling distance=4.2cm,
  every node/.style={draw,thick,rounded corners=3pt,align=center,font=\small,inner sep=5pt},
  edge from parent/.style={draw,thick,-{Stealth}}]
  \node[fill=kcblue!20] {\textbf{My garden picture}\\(main script)}
    child {node[fill=kcpurple!25] {\textbf{grass}\\green strip at the bottom}}
    child {node[fill=kcpurple!25] {\textbf{flower (size)}\\stamp one flower}}
    child {node[fill=kcpurple!25] {\textbf{repeat 12}\\random places,\\random sizes}};
\end{tikzpicture}
\end{center}

In Scratch, \menu{My Blocks} $\rightarrow$ \menu{Make a Block} lets you make a pink block with your own name. You put its steps under a \textbf{define} hat. An \textbf{input} is a blank you fill in each time you use the block, like \enquote{size}.

\subsection{Messages}
\textbf{broadcast} sends a message to \textbf{every} sprite, like an announcement on a loudspeaker. Only scripts that start with \textbf{when I receive} \textbf{that exact message} run.

\begin{center}
\begin{tikzpicture}[>={Stealth},
  spr/.style={draw,thick,rounded corners=3pt,align=center,font=\small,inner sep=5pt,minimum width=3.4cm}]
  \node[spr,fill=kcorange!20] (cat) at (0,0) {\faCat\ \textbf{Presenter}\\broadcast \enquote{start show}};
  \node[spr,fill=kcgreen!20] (d) at (7,0.9) {\textbf{Dancer}\\when I receive \enquote{start show}\\ \textbf{✓ starts dancing}};
  \node[spr,fill=kcblack!8] (g) at (7,-1.1) {\textbf{Ghost}\\when I receive \enquote{show start}\\ \textbf{✗ nothing happens}};
  \draw[->,thick] (cat.east) -- node[above,font=\small]{\faBullhorn} (d.west);
  \draw[->,thick,dashed] (cat.east) -- (g.west);
\end{tikzpicture}
\end{center}

\begin{exercise}[Shape Machine]
Make your own block that can draw \textbf{any} shape.
\begin{enumerate}[leftmargin=1.4em,itemsep=1pt]
  \item New project. Add the \menu{Pen} extension. Set the cat's \menu{Size} to 40.
  \item \menu{My Blocks} $\rightarrow$ \menu{Make a Block}. Type \texttt{polygon}. Click \menu{Add an input (number or text)}, type \texttt{sides}. Add another input: \texttt{length}. Click \menu{OK}.
  \item Build the define script. Drag the \texttt{sides} and \texttt{length} ovals \textbf{out of the define hat} into the slots.
  \item Build the green-flag script with your new pink block. Click the flag!
  \item Add a pentagon: one more block, \enquote{polygon 5 70}.
\end{enumerate}
\begin{center}
\begin{scratch}[scale=0.78]
\initmoreblocks{define \namemoreblocks{polygon \ovalmoreblocks{sides} \ovalmoreblocks{length}}}
\blockpen{pen down}
\blockrepeat{repeat \ovalmoreblocks{sides}}{
  \blockmove{move \ovalmoreblocks{length} steps}
  \blockmove{turn \turnright{} \ovaloperator{\ovalnum{360} / \ovalmoreblocks{sides}} degrees}
}
\blockpen{pen up}
\end{scratch}
\hspace{0.8cm}
\begin{scratch}[scale=0.78]
\blockinit{when \greenflag clicked}
\blockpen{erase all}
\blockmove{go to x: \ovalnum{0} y: \ovalnum{0}}
\blockmove{point in direction \ovalnum{90}}
\blockmoreblocks{polygon \ovalnum{3} \ovalnum{120}}
\blockmoreblocks{polygon \ovalnum{4} \ovalnum{100}}
\blockmoreblocks{polygon \ovalnum{6} \ovalnum{60}}
\blockmoreblocks{polygon \ovalnum{8} \ovalnum{40}}
\end{scratch}
\end{center}
How many blocks would you need to draw these 4 shapes \textbf{without} your My Block? \rule{2cm}{0.4pt}
\end{exercise}

\begin{exercise}[Stamp Garden]
\begin{enumerate}[leftmargin=1.4em,itemsep=1pt]
  \item New project with the \menu{Pen} extension. Delete the cat. Hover over \menu{Choose a Sprite} and click the \menu{Paint} brush. Draw a flower in the middle of the canvas (circles for the center and petals). Name the sprite \texttt{Flower}.
  \item Make two My Blocks: \texttt{flower} with an input \texttt{size}, and \texttt{grass} with no input.
  \item Build the three scripts. For the grass color, click the color circle and set \textbf{Color 45, Saturation 100, Brightness 62} (see the color-safe paint box).
\end{enumerate}
\begin{center}
\begin{scratch}[scale=0.72]
\initmoreblocks{define \namemoreblocks{flower \ovalmoreblocks{size}}}
\blocklook{set size to \ovalmoreblocks{size} \%}
\blocklook{change \selectmenu{color} effect by \ovalnum{25}}
\blockpen{stamp}
\end{scratch}
\hspace{0.4cm}
\begin{scratch}[scale=0.72]
\initmoreblocks{define \namemoreblocks{grass}}
\blockpen{set pen color to \pencolor{kcgreen}}
\blockpen{set pen size to \ovalnum{60}}
\blockmove{go to x: \ovalnum{-240} y: \ovalnum{-150}}
\blockpen{pen down}
\blockmove{go to x: \ovalnum{240} y: \ovalnum{-150}}
\blockpen{pen up}
\end{scratch}

\medskip
\begin{scratch}[scale=0.72]
\blockinit{when \greenflag clicked}
\blockpen{erase all}
\blocklook{clear graphic effects}
\blocklook{show}
\blockmoreblocks{grass}
\blockrepeat{repeat \ovalnum{12}}{
  \blockmove{go to x: \ovaloperator{pick random \ovalnum{-220} to \ovalnum{220}} y: \ovaloperator{pick random \ovalnum{-140} to \ovalnum{120}}}
  \blockmoreblocks{flower \ovaloperator{pick random \ovalnum{30} to \ovalnum{100}}}
}
\blocklook{hide}
\end{scratch}
\end{center}
\textbf{Done when:} each click on the flag grows a new garden: a strip of grass and 12 flowers of different sizes and colors in surprise places.
\end{exercise}

\begin{exercise}[Show Time]
Two sprites put on a show by sending messages.
\begin{enumerate}[leftmargin=1.4em,itemsep=1pt]
  \item New project. Keep the cat (the \textbf{Presenter}). Add a \textbf{Dancer}: any sprite with several costumes, like Ballerina.
  \item To make a message: in a \enquote{broadcast} block, click the dropdown $\rightarrow$ \menu{New message}, type \texttt{start show}. Make another called \texttt{finale}.
  \item Build the scripts. Click the flag and enjoy the show!
\end{enumerate}
\begin{center}
\begin{tabular}{@{}c@{\hspace{0.8cm}}c@{}}
\begin{scratch}[scale=0.72]
\blockinit{when \greenflag clicked}
\blockmove{go to x: \ovalnum{-150} y: \ovalnum{-50}}
\blocklook{say \ovalnum{Welcome to the show!} for \ovalnum{2} seconds}
\blockevent{broadcast \selectmenu{start show}}
\end{scratch}
&
\begin{scratch}[scale=0.72]
\blockinit{when \greenflag clicked}
\blocklook{hide}
\end{scratch}
\\[6pt]
\begin{scratch}[scale=0.72]
\blockinit{when I receive \selectmenu{finale}}
\blocklook{say \ovalnum{Thank you for watching!} for \ovalnum{2} seconds}
\end{scratch}
&
\begin{scratch}[scale=0.72]
\blockinit{when I receive \selectmenu{start show}}
\blockmove{go to x: \ovalnum{80} y: \ovalnum{-30}}
\blocklook{show}
\blockrepeat{repeat \ovalnum{8}}{
  \blocklook{next costume}
  \blockcontrol{wait \ovalnum{0.3} seconds}
}
\blockevent{broadcast \selectmenu{finale}}
\end{scratch}
\\[4pt]
{\small\textbf{Presenter} (cat)} & {\small\textbf{Dancer}}
\end{tabular}
\end{center}
\textbf{Bug hunt!} Ada's dancer starts with \enquote{when I receive \textbf{show start}}. Try it. What happens? Why? How do you fix it?
\answerlines{2}
\end{exercise}

\begin{safety}[Color-safe paint box]
Some people see colors differently. These colors stay different for almost everyone. In Scratch, click a color circle and type the three numbers. And never use color as the only clue: add a shape, a pattern or a word too, and keep text big.

\begin{center}
\renewcommand{\arraystretch}{1.25}
\begin{tabular}{@{}llccc@{\hspace{1.2cm}}llccc@{}}
 & \textbf{Color} & \textbf{Col.} & \textbf{Sat.} & \textbf{Bri.} & & \textbf{Color} & \textbf{Col.} & \textbf{Sat.} & \textbf{Bri.} \\
\swatch{kcorange} & Orange & 12 & 100 & 90 & \swatch{kcblue} & Blue & 56 & 100 & 70 \\
\swatch{kcsky} & Sky blue & 56 & 63 & 91 & \swatch{kcvermillion} & Vermillion & 7 & 100 & 84 \\
\swatch{kcgreen} & Bluish green & 45 & 100 & 62 & \swatch{kcpurple} & Reddish purple & 91 & 41 & 80 \\
\swatch{kcyellow} & Yellow & 16 & 72 & 94 & \swatch{kcblack} & Black & 0 & 0 & 0 \\
\end{tabular}
\end{center}
\end{safety}

\begin{vocab}
  \vocabword{decomposition}{Breaking a big problem into smaller parts.}
  \vocabword{My Block}{A block you make yourself, with a name you choose.}
  \vocabword{define}{The hat block that holds the steps of your My Block.}
  \vocabword{input}{A blank in a block that you fill in each time, like \enquote{size}.}
  \vocabword{broadcast}{Send a message to every sprite.}
  \vocabword{when I receive}{A hat block that runs when its exact message is sent.}
  \vocabword{accessible}{Easy for everyone to use, including people who see or hear differently.}
\end{vocab}

\begin{reflection}
Why is \enquote{polygon 6 60} easier to read than the blocks inside it?
\answerlines{2}
What did you do to make your project easier for everyone to see?
\answerlines{1}
\end{reflection}
